\documentclass[11pt,a4paper]{article}

\usepackage[T1]{fontenc}
\usepackage[utf8]{inputenc}
\usepackage[brazilian]{babel}
\usepackage{lmodern}
\usepackage[varqu,varl]{zi4}
\usepackage{microtype}
\usepackage[a4paper,margin=2.5cm,headheight=14pt]{geometry}
\usepackage{amsmath,amssymb}
\usepackage[color]{coqdoc}
\usepackage{xcolor}
\usepackage[most]{tcolorbox}
\usepackage{enumitem}
\usepackage{fancyhdr}
\usepackage{titlesec}
\usepackage{fontawesome5}

% ------------------------------------------------------------------
% Cores
% ------------------------------------------------------------------
\definecolor{unbblue}{RGB}{0,63,135}
\definecolor{unbgreen}{RGB}{0,122,61}
\definecolor{codeback}{RGB}{248,249,252}
\definecolor{codeframe}{RGB}{190,200,220}

\hypersetup{colorlinks=true,linkcolor=unbblue,urlcolor=unbblue,citecolor=unbblue,
  pdftitle={LC1 2026/2 -- Projeto: Árvores Binárias de Busca},
  pdfauthor={Flávio L. C. de Moura}}

% Liga as definições nos blocos de código ao repositório. Os arquivos
% coqsrcbase.tex e coqsrclinks.tex são gerados por `make doc`.
\input{coqsrcbase}
\makeatletter
\newcommand{\coqsrcline}[3]{\@namedef{coqsrc@#1}{#2\#L#3}}
\newcommand{\coqsrclink}[2]{\@ifundefined{coqsrc@#1}{#2}{\href{\coqsrcbase\@nameuse{coqsrc@#1}}{#2}}}
\makeatother
\input{coqsrclinks}

% ------------------------------------------------------------------
% Títulos e cabeçalhos
% ------------------------------------------------------------------
\titleformat{\section}{\Large\bfseries\color{unbblue}}{\thesection}{0.6em}{}
  [{\color{unbblue!40}\titlerule[0.8pt]}]
\titleformat{\subsection}{\large\bfseries\color{unbblue!85!black}}{\thesubsection}{0.6em}{}

\pagestyle{fancy}
\fancyhf{}
\fancyhead[L]{\small\color{unbblue}Lógica Computacional 1 -- 2026/2}
\fancyhead[R]{\small\color{unbblue}Projeto: árvores binárias de busca}
\fancyfoot[C]{\small\thepage}
\renewcommand{\headrulewidth}{0.4pt}
\renewcommand{\headrule}{\hbox to\headwidth{\color{unbblue!40}\leaders\hrule height \headrulewidth\hfill}}

\setlength{\parskip}{0.5em}
\setlength{\parindent}{0pt}

% ------------------------------------------------------------------
% Código Rocq gerado pelo rocq doc em caixas tcolorbox
% ------------------------------------------------------------------
% tools/tidy_coqdoc.py passa a chave da primeira definição do bloco como
% argumento opcional de coqdoccode; o título da caixa aponta para ela no
% repositório.
\newtcolorbox{rocqbox}[1]{enhanced,
  colback=codeback,colframe=codeframe,boxrule=0.5pt,arc=2mm,
  leftrule=3pt,left=8pt,right=8pt,top=6pt,bottom=6pt,
  borderline west={3pt}{0pt}{unbblue},
  before skip=0.8em,after skip=0.8em,
  attach boxed title to top right={yshift=-2mm,xshift=-3mm},
  boxed title style={colback=unbblue,colframe=unbblue,arc=1mm,boxrule=0pt,
    left=3pt,right=3pt,top=1pt,bottom=1pt},
  coltitle=white,fonttitle=\scriptsize\sffamily\bfseries,
  fontupper=\small,
  title={\coqsrclink{#1}{\textcolor{white}{Rocq}}}}

\RenewDocumentEnvironment{coqdoccode}{O{}}
  {\begin{rocqbox}{#1}}
  {\end{rocqbox}}

% Links para a biblioteca padrão do Rocq
\renewcommand{\coqexternalref}[3]{\href{#2.html}{#3}}

% Caixas de destaque para o texto do enunciado
\newtcolorbox{objetivo}{enhanced,colback=unbgreen!6,colframe=unbgreen,
  boxrule=0.6pt,arc=2mm,left=10pt,right=10pt,
  fonttitle=\bfseries,title={\faBullseye\quad Objetivo do projeto}}
\newtcolorbox{atencao}{enhanced,colback=orange!6,colframe=orange!80!black,
  boxrule=0.6pt,arc=2mm,left=10pt,right=10pt,
  fonttitle=\bfseries,title={\faExclamationTriangle\quad Atenção}}

\newcommand{\rocq}{\textsf{Rocq}}
\newcommand{\prazo}{defesa oral, com data a agendar}

\begin{document}

% ------------------------------------------------------------------
% Cabeçalho
% ------------------------------------------------------------------
\thispagestyle{empty}
\begin{tcolorbox}[enhanced,colback=unbblue,colframe=unbblue,arc=0mm,
  left=12pt,right=12pt,top=10pt,bottom=10pt,coltext=white]
  {\small Universidade de Brasília \hfill Departamento de Ciência da Computação}\\[0.6em]
  {\LARGE\bfseries Lógica Computacional 1 -- 2026/2}\\[0.4em]
  {\Large Projeto: Formalização de árvores binárias de busca (BST)}\\[0.6em]
  {\small Prof. Flávio L. C. de Moura \hfill Entrega: \prazo}
\end{tcolorbox}


\section{Introdução}

Uma \emph{árvore binária de busca} (BST, Binary Search Tree) é uma árvore binária em que cada nó
armazena uma chave (possivelmente acompanhada de um valor) de forma que,
para todo nó com chave $k$, as chaves da subárvore à esquerda são menores
do que $k$ e as chaves da subárvore à direita são maiores do que $k$. Esse
invariante permite implementar as operações de \emph{busca},
\emph{inserção} e \emph{remoção} percorrendo um único caminho da raiz até
uma folha, isto é, em tempo proporcional à altura da árvore: $O(\log n)$
quando a árvore está balanceada, mas $\Theta(n)$ no pior caso. As árvores
binárias de busca são a base de estruturas de dados como as árvores AVL,
as árvores rubro-negras e as \emph{treaps}, e são utilizadas na
implementação de dicionários e conjuntos em muitas linguagens de
programação. Uma apresentação detalhada pode ser encontrada no Capítulo~12
do livro de Cormen, Leiserson, Rivest e Stein~\cite{cormen2022}.

As árvores binárias de busca foram descobertas de forma independente por
diversos pesquisadores no final da década de 1950, em um momento em que se
buscavam formas eficientes de armazenar tabelas de símbolos e dicionários
na memória dos primeiros computadores. Os primeiros trabalhos publicados
são os de P.~F.~Windley~\cite{windley1960} e de A.~D.~Booth e
A.~J.~T.~Colin~\cite{booth1960}, ambos de 1960. Em 1962, T.~N.~Hibbard
publicou um estudo das propriedades combinatórias dessas
árvores~\cite{hibbard1962}, no qual descreve o algoritmo de remoção que
estudaremos neste projeto: para remover a chave de um nó com duas
subárvores não vazias, ela é substituída pela menor chave da subárvore à
direita (o seu \emph{sucessor}), que é então removida dessa subárvore.

\begin{tcolorbox}[enhanced,breakable,colback=unbblue!4,colframe=unbblue!60,
  boxrule=0.6pt,arc=2mm,left=10pt,right=10pt,fonttitle=\bfseries,
  title={\faLightbulb[regular]\quad Curiosidades}]
  \begin{itemize}[leftmargin=1.2em,itemsep=0.3em]
  \item Se $n$ chaves distintas são inseridas em uma ordem aleatória, a
    altura esperada da árvore resultante é $O(\log n)$. Por outro lado,
    se as chaves são inseridas em ordem crescente, a árvore degenera em
    uma lista e as operações passam a custar $\Theta(n)$. É por isso que,
    na prática, utilizam-se variantes \emph{balanceadas}, que realizam
    rotações para manter a altura logarítmica.
  \item A remoção é muito mais difícil de analisar do que a inserção.
    Hibbard mostrou que uma remoção preserva a ``aleatoriedade'' de uma
    árvore aleatória, mas sequências alternadas de inserções e remoções
    modificam a distribuição das árvores de forma sutil. A análise de um
    caso aparentemente simples, com árvores de apenas três nós, rendeu a
    Jonassen e Knuth um artigo inteiro, com o sugestivo título \emph{A
      trivial algorithm whose analysis isn't}~\cite{jonassen1978}.
  \item Como a remoção de Hibbard sempre utiliza o sucessor, e nunca o
    predecessor, ela é assimétrica. Experimentos de
    Eppinger~\cite{eppinger1983} indicaram que longas sequências de
    inserções e remoções aleatórias tendem a deixar a árvore
    \emph{menos} balanceada do que uma árvore construída apenas por
    inserções.
  \end{itemize}
\end{tcolorbox}

\subsection*{O projeto}

Neste projeto, as árvores binárias de busca serão estudadas do ponto de
vista da \emph{verificação formal}: em vez de testar as operações em
alguns exemplos, queremos construir \emph{provas matemáticas}, verificadas
mecanicamente pelo assistente de provas \rocq{}~(versão 9.1.1), de que
elas estão corretas para \emph{todas} as árvores possíveis. A formalização
apresentada a seguir foi adaptada da referência~\cite{appel} (capítulo
\emph{SearchTree} do livro \emph{Verified Functional Algorithms}, volume~3
da série \emph{Software Foundations}) e estendida com a operação de remoção.

\begin{objetivo}
  O objetivo deste projeto é \textbf{provar formalmente, no assistente de
    provas \rocq, a correção das operações de busca, inserção e remoção em
    árvores binárias de busca}, isto é, que:
  \begin{enumerate}[label=(\roman*),nosep]
  \item as funções \texttt{insert} e \texttt{delete} preservam o
    invariante \texttt{BST};
  \item as funções de busca \texttt{bound} e \texttt{lookup} se comportam
    como esperado após inserções e remoções.
  \end{enumerate}
\end{objetivo}

\section{A formalização}

Todo o código-fonte deste documento está disponível no
\href{https://github.com/flaviodemoura/bst-search-insert-delete}{repositório
do projeto}. Em cada caixa de código, o título \rocq{} é um \emph{link}
para o trecho correspondente no repositório.

A seguir, apresentamos as definições das árvores binárias, do invariante
\texttt{BST} e das operações de busca, inserção e remoção, bem como as
questões do projeto.

\input{bst.v.tex}

\section{Instruções}

\begin{itemize}[leftmargin=*]
\item O projeto pode ser feito \textbf{individualmente ou em grupos de até
    4 alunos}.
\item A tarefa consiste em resolver as 9 questões propostas. 
\item É permitido (e recomendado) definir lemas auxiliares e utilizar os
  resultados da biblioteca padrão do \rocq. Algumas dicas:
  \begin{itemize}[nosep]
  \item a tática \texttt{bdall} resolve automaticamente muitos dos casos
    gerados pelas comparações \texttt{<?} entre chaves, e a tática
    \texttt{bdestruct} permite fazer essa análise de casos manualmente;
  \item nas questões sobre a remoção, a indução deve ser feita com cuidado:
    observe sobre qual argumento a função \texttt{remove\_min} é recursiva,
    e quais argumentos devem ser generalizados antes da indução;
  \item o lema \texttt{ForallT\_impl} é útil para obter, por exemplo,
    \texttt{ForallT (fun y \_ => y < m) l} a partir de
    \texttt{ForallT (fun y \_ => y < x) l} e $x < m$.
  \end{itemize}
\end{itemize}

\begin{atencao}
  As provas devem compilar sem erros no \rocq{} e não podem conter
  os comandos \texttt{Admitted}, \texttt{admit} ou \texttt{Axiom}.
\end{atencao}

\section{Uso de IA e entrega}

\begin{itemize}[leftmargin=*]
\item O projeto pode ser desenvolvido com o apoio de ferramentas de
  inteligência artificial. No entanto, \textbf{a compreensão da resolução faz
    parte do projeto}: não basta apresentar uma prova que compila, é
  necessário entendê-la. Cada integrante do grupo deve ser capaz de
  explicar a estratégia geral da prova, o papel de cada lema auxiliar e o
  efeito das táticas utilizadas, bem como de responder a perguntas sobre
  possíveis modificações da prova ou do enunciado.
\item A entrega será feita por meio de uma \textbf{apresentação e defesa
    oral} de cada grupo, em data a ser agendada ao longo das próximas aulas.
\item Os grupos devem estar preparados para \textbf{explicar as provas em
    linguagem natural}, isto é, como uma prova matemática usual, sem
  recorrer às táticas do \rocq: por exemplo, sobre o que é feita a indução,
  quais são os casos e por que os lemas auxiliares são necessários.
  \textbf{Essa explicação fará parte da avaliação.}
\item Durante a defesa, perguntas poderão ser dirigidas a qualquer
  integrante do grupo, e a avaliação levará em conta tanto a prova
  apresentada quanto a compreensão demonstrada por cada um.
\end{itemize}

\begin{tcolorbox}[enhanced,colback=unbblue!4,colframe=unbblue!60,
  boxrule=0.6pt,arc=2mm,left=10pt,right=10pt,fonttitle=\bfseries,
  title={\faRobot\quad Sobre o uso de IA}]
  Use a IA como uma ferramenta de apoio ao aprendizado: peça explicações,
  explore alternativas e questione as respostas obtidas. Uma prova que o grupo
  não consegue explicar não será considerada uma solução do projeto.
\end{tcolorbox}

\section{Outros temas}

O tema deste projeto (a formalização de árvores binárias de busca) é uma
\textbf{sugestão}. Os grupos podem desenvolver outros temas, \textbf{desde
  que o tema seja acordado previamente com o professor}, antes do início do
seu desenvolvimento. A utilização de outros assistentes de prova (Lean,
Agda, PVS, Isabelle/HOL, etc) também é possível. As demais instruções
deste documento, em particular as relativas ao uso de IA e à defesa oral,
também se aplicam a esses temas.

{\small
\begin{thebibliography}{9}
\setlength{\itemsep}{1pt}
\bibitem{cormen2022} T.~H. Cormen, C.~E. Leiserson, R.~L. Rivest e C.~Stein.
  \emph{Introduction to Algorithms}. 4.\textsuperscript{a} edição, MIT Press,
  2022.
\bibitem{windley1960} P.~F. Windley. Trees, forests and rearranging.
  \emph{The Computer Journal}, 3(2):84--88, 1960.
\bibitem{booth1960} A.~D. Booth e A.~J.~T. Colin. On the efficiency of a
  new method of dictionary construction. \emph{Information and Control},
  3(4):327--334, 1960.
\bibitem{hibbard1962} T.~N. Hibbard. Some combinatorial properties of
  certain trees with applications to searching and sorting. \emph{Journal
    of the ACM}, 9(1):13--28, 1962.
\bibitem{jonassen1978} A.~T. Jonassen e D.~E. Knuth. A trivial algorithm
  whose analysis isn't. \emph{Journal of Computer and System Sciences},
  16(3):301--322, 1978.
\bibitem{eppinger1983} J.~L. Eppinger. An empirical study of insertion and
  deletion in binary search trees. \emph{Communications of the ACM},
  26(9):663--669, 1983.
\bibitem{appel} A.~W. Appel. \emph{Verified Functional Algorithms}.
  Volume~3 de \emph{Software Foundations}, B.~C. Pierce (ed.).
  \url{https://softwarefoundations.cis.upenn.edu/}.
\end{thebibliography}
}

\end{document}
